%% 7장 조합론과 생성함수
\chapter{조합론과 생성함수}
\label{ch:combinatorics}

\providecommand\qbinom[2]{\genfrac{[}{]}{0pt}{}{#1}{#2}_{q}}

\abstract{순위 검정의 정확한 p값은 확률을 적분해서 얻는 것이 아니라 경우의 수를 세어서 얻는다. 귀무가설 아래에서 관측값들의 순서가 모두 같은 확률로 나타난다면, p값은 관측한 통계량만큼 극단적인 배치의 수를 전체 배치의 수로 나눈 몫이기 때문이다. 이 장은 그 수를 세는 수학과 알고리즘을 다룬다. 먼저 이항계수와 초기하 확률을 넘치지 않게 계산하는 방법을 보고, 수열을 멱급수의 계수로 담는 생성함수를 형식적 멱급수의 환으로 정식화한다. 이어서 Gauss가 도입한 $q$-이항계수가 Mann--Whitney 통계량의 분포 그 자체임을 증명하고, 부호 순위 통계량은 부분집합의 합으로, Kendall의 $\tau$는 순열의 역순의 수로 바꾸어 각각의 생성함수와 동적 계획 알고리즘을 얻는다. 동점이 있을 때의 두 변수 생성함수, 행의 순열을 부분집합으로 세는 Page 검정을 거쳐, 마지막으로 두 표본 Kolmogorov--Smirnov 검정의 분포를 격자 경로의 반사 원리와 흡수 경계가 있는 동적 계획으로 구한다.}

두 가지 처리를 받은 실험 단위가 각각 $3$개와 $4$개 있고, 처리 효과를 측정한 일곱 값을 크기 순으로 늘어놓았더니 첫째 처리의 값이 $1$번째, $2$번째, $4$번째 자리에 왔다고 하자. 첫째 처리가 값을 낮추는 효과가 있다는 증거로 볼 수 있을까? 두 처리에 차이가 없다면 일곱 값은 한 분포에서 나온 셈이므로, 첫째 처리의 세 값이 일곱 자리 가운데 어느 세 자리를 차지할지는 $\binom73=35$가지가 모두 같은 확률로 일어난다. 관측한 배치보다 첫째 처리가 더 앞쪽에 몰린 배치는 자리 $\{1,2,3\}$ 하나뿐이고 관측한 배치 자체도 있으므로, ``이만큼 또는 이보다 더 앞에 몰리는'' 확률은 $2/35\approx0.057$이다. 이 계산에는 분포를 가정할 필요도, 적분할 필요도 없었고 경우를 센 것이 전부였다. 다만 ``더 앞쪽에 몰렸다''를 재는 통계량을 정해야 하고, 표본이 커지면 일일이 나열할 수 없으므로 세는 방법이 필요하다. 이 장은 바로 그것을 다룬다.

\section{세기의 기본}
\label{sec:counting}

\subsection{순열과 이항계수}

서로 다른 $n$개의 대상을 한 줄로 늘어놓는 방법의 수는 첫 자리에 $n$가지, 둘째 자리에 남은 $n-1$가지를 놓는 식으로 곱해 $n!=n(n-1)\cdots1$이다. 이처럼 선택이 단계별로 이루어지고 각 단계의 선택지 수가 앞 단계의 선택에 상관없이 일정하면 전체 경우의 수는 단계별 수의 곱이 된다는 원리를 \term{곱의 법칙}{rule of product}이라 한다. $n$개 가운데 $k$개를 골라 순서 있게 늘어놓는 수는 같은 원리로 $n(n-1)\cdots(n-k+1)$이고, 고른 $k$개의 순서를 잊으면 같은 집합이 $k!$번씩 세어지므로 순서 없이 $k$개를 고르는 수는
\[
\binom nk=\frac{n(n-1)\cdots(n-k+1)}{k!}=\frac{n!}{k!\,(n-k)!}
\]
이다. 이것이 \term{이항계수}{binomial coefficient}이며, $k<0$이나 $k>n$이면 $0$으로 둔다.

이항계수의 기본 성질은 모두 ``무엇을 세는가''에서 바로 나온다. 고른 $k$개와 고르지 않은 $n-k$개는 서로를 결정하므로 $\binom nk=\binom n{n-k}$이다. 원소 $n$을 고르는 경우와 고르지 않는 경우로 나누면 Pascal의 공식
\begin{equation}
\label{eq:pascal}
\binom nk=\binom{n-1}{k-1}+\binom{n-1}k
\end{equation}
을 얻는다. $(1+t)^n$을 $n$개의 인수 $(1+t)$의 곱으로 보고 전개하면 각 인수에서 $1$과 $t$ 중 하나를 고르게 되는데, $t$를 고른 인수가 $k$개인 항이 $t^k$를 주므로 \term{이항정리}{binomial theorem} $(1+t)^n=\sum_k\binom nkt^k$가 성립한다. 같은 방식으로 $(1+t)^{m+n}=(1+t)^m(1+t)^n$의 $t^k$ 계수를 비교하면 Vandermonde의 항등식
\begin{equation}
\label{eq:vandermonde}
\binom{m+n}k=\sum_{i}\binom mi\binom n{k-i}
\end{equation}
을 얻는데, 이것은 $m+n$개에서 $k$개를 고를 때 앞쪽 $m$개에서 몇 개를 고르는지로 경우를 나눈 것이다. 생성함수의 곱이 계수들의 합성곱을 준다는 사실을 처음 만나는 자리이기도 하다.

이항계수는 빨리 자란다. Stirling 공식(3장)에서 $\binom{2n}n\sim4^n/\sqrt{\pi n}$이므로 $\binom{2000}{1000}$은 약 $2.0\times10^{600}$이고, 이 값은 \code{double}이 나타낼 수 있는 가장 큰 수 약 $1.8\times10^{308}$을 훌쩍 넘는다. 표본 크기가 수천인 검정에서 경우의 수를 그대로 세면 곧바로 넘침을 만난다는 뜻이므로, 이 장의 알고리즘들은 정확한 정수로 셀 수 있는 범위와 그 너머에서 확률로 바꾸어 쌓는 방법을 함께 갖춘다.

\subsection{이항계수를 넘치지 않게 계산하기}

$\binom nk$를 정수로 정확히 구할 때 $n!/(k!(n-k)!)$을 그대로 계산하면 분자가 답보다 훨씬 먼저 넘친다. 곱셈 공식 $\binom ni=\binom n{i-1}\cdot(n-i+1)/i$를 $i=1,\dots,k$로 되풀이하면 중간값이 모두 이항계수이므로 낫지만, 곱 $\binom n{i-1}(n-i+1)$이 답의 최대 $n$배까지 커지는 문제가 남는다. 2장의 최대공약수를 쓰면 이 곱셈조차 답을 넘지 않게 할 수 있다.

\begin{proposition}
\label{prop:binom-gcd}
$c=\binom n{i-1}$, $g=\gcd(c,i)$라 하자. 그러면 $i/g$는 $n-i+1$을 나누고
\[
\binom ni=\frac cg\cdot\frac{n-i+1}{i/g}
\]
이다. 우변의 두 인수는 모두 $\binom ni$ 이하의 양의 정수이다.
\end{proposition}

\begin{proof}
$\binom ni=c(n-i+1)/i$가 정수이므로 $i$는 $c(n-i+1)$을 나눈다. 양변을 $g$로 나누면 $i/g$가 $(c/g)(n-i+1)$을 나누는데, $\gcd(c/g,i/g)=\gcd(c,i)/g=1$이므로 2장의 Euclid 보조정리에서 $i/g$가 $n-i+1$을 나눈다. 따라서 두 인수 $c/g$와 $(n-i+1)/(i/g)$는 양의 정수이고, 그 곱이 $\binom ni$이므로 각 인수는 $\binom ni$를 넘지 않는다.
\end{proof}

\cref{prop:binom-gcd}에 따라 계산의 모든 중간값이 최종 답 이하이므로, 답이 들어가는 정수형이라면 계산 도중에 넘치는 일이 없다. 대가는 단계마다 최대공약수 하나인데, $i\le k$이므로 Euclid 알고리즘은 $O(\log k)$번의 나눗셈으로 끝난다.

\begin{instardust}[\code{csrc/htest/rank\_exact.hpp}]
순위 검정의 정확 분포는 경우의 수가 $2^{124}$ 아래일 때 $128$비트 정수로 센다. 이항계수 \code{binom\_u128}은 \cref{prop:binom-gcd}의 순서 그대로 공통 인수를 먼저 덜어 낸 뒤 곱한다. 주석에 적힌 대로 어떤 중간값도 결과를 넘지 않으므로, 결과가 $2^{126}$ 아래라는 호출자의 약속만으로 넘침이 없다는 것이 보장된다.
\end{instardust}

\subsection{초기하 확률}

$R$개의 자리 가운데 $r$개에 표시가 되어 있고, 그 가운데 $t$개를 무작위로 뽑을 때 표시된 자리가 $k$개 들어 있을 확률은 뽑은 $t$개 중 표시된 $k$개와 표시되지 않은 $t-k$개를 각각 고르는 수의 곱을 전체로 나눈
\begin{equation}
\label{eq:hypergeom}
h(k)=\binom rk\binom{R-r}{t-k}\Bigm/\binom Rt
\end{equation}
이다. 이 분포를 \term{초기하 분포}{hypergeometric distribution}라 하고, Vandermonde 항등식 \cref{eq:vandermonde}이 곧 $\sum_kh(k)=1$이다. 셋째 인수를 바꾸어 $\binom tk\binom{R-t}{r-k}/\binom Rr$로 써도 같은 값인데, 뽑힌 자리와 표시된 자리의 역할을 바꾸어 센 것이기 때문이다.

세 이항계수는 각각 넘칠 수 있지만 확률 자체는 $1$ 이하이므로, 계수를 따로 구하지 않고 이웃한 항의 비
\[
\frac{h(k+1)}{h(k)}=\frac{(t-k)(r-k)}{(k+1)(R-t-r+k+1)}
\]
로 항들을 차례로 만드는 것이 좋다. 비가 $1$보다 큰 동안 $h$는 증가하고 그 뒤로 감소하므로 분포는 \term{단봉}{unimodal}이며, 비가 $1$ 이상이 되는 조건을 정리하면 최빈값이 $\lfloor(t+1)(r+1)/(R+2)\rfloor$임을 얻는다. 그래서 최빈값에 임시로 $1$을 놓고 양쪽으로 비를 곱해 나가면 모든 값이 $1$ 이하로 줄어들기만 하므로 넘침이 없고, 아래로 넘쳐 $0$이 된 항에서 멈출 수 있다. 끝으로 전체의 합으로 나누면 정규화된 확률을 얻는다. 합이 모두 양수의 덧셈이라는 점도 중요한데, 양수만 더하는 합은 상대오차가 덧셈 횟수에 비례하는 만큼만 쌓인다는 사실을 10장에서 증명한다.

\begin{instardust}[\code{csrc/htest/ks2.hpp}]
동점이 있는 두 표본 Kolmogorov--Smirnov 검정은 같은 값들의 묶음마다 그 묶음이 어느 표본의 값을 몇 개 갖는지를 초기하 확률로 다룬다(\cref{sec:lattice}). \code{hypergeometric} 함수는 바로 위의 방법대로 최빈값에서 비를 곱하고 나누어 항을 만들고, 비 자체는 분자와 분모를 각각 정확한 곱(10장의 TwoProd)으로 만든 double-double의 몫이다.
\end{instardust}

\section{생성함수}
\label{sec:gf}

\subsection{수열을 담는 멱급수}

수열 $a_0,a_1,a_2,\dots$를 멱급수 $A(q)=\sum_ka_kq^k$의 계수로 담은 것을 그 수열의 \term{생성함수}{generating function}라 한다. 이렇게 하는 첫째 이유는 곱셈에 있다. 두 생성함수의 곱은
\[
A(q)B(q)=\sum_s\Bigl(\sum_{k}a_kb_{s-k}\Bigr)q^s
\]
이므로 곱의 계수는 두 수열의 \term{합성곱}{convolution}이다. 경우의 수로 읽으면, 크기가 $k$인 대상이 $a_k$가지이고 크기가 $l$인 대상이 $b_l$가지일 때 둘을 하나씩 짝지어 크기를 더한 것이 $s$인 경우의 수가 곱의 $q^s$ 계수이다.

확률로 읽으면 더 쓸모가 크다. 음이 아닌 정수 값을 갖는 확률변수 $X$에 대해 $G_X(q)=E[q^X]=\sum_kP(X=k)q^k$를 $X$의 \term{확률생성함수}{probability generating function}라 하는데, $X$와 $Y$가 독립이면 $E[q^{X+Y}]=E[q^X]E[q^Y]$이므로 독립인 확률변수의 합의 분포는 확률생성함수의 곱으로 얻는다. 주사위 두 개의 합이라면 $\bigl((q+q^2+\cdots+q^6)/6\bigr)^2$을 전개해 $q^7$의 계수 $6/36$을 읽으면 된다. 이 장의 순위 통계량들은 거의 모두 독립이거나 독립에 가까운 조각들의 합이므로, 분포를 구하는 일은 곱 꼴의 생성함수를 찾고 그 계수를 효율적으로 전개하는 일이 된다.

\subsection{형식적 멱급수의 환}

생성함수를 다룰 때 $q$에 수를 넣거나 수렴을 따질 필요는 없다. 계수의 수열 자체를 대상으로 보고 덧셈과 곱셈을 위의 규칙으로 정의하면 된다.

\begin{definition}[형식적 멱급수]
\label{def:formal-power-series}
가환환 $K$의 원소로 이루어진 수열 $(a_k)_{k\ge0}$을 $\sum_ka_kq^k$로 적고 \term{형식적 멱급수}{formal power series}라 한다. 덧셈은 계수별로, 곱셈은 합성곱으로 정의하며, 이 연산으로 이루는 환을 $K[[q]]$로 적는다. $q^s$의 계수를 $[q^s]A$로 쓴다.
\end{definition}

곱의 각 계수 $\sum_{k=0}^sa_kb_{s-k}$는 유한합이므로 정의에는 아무런 수렴 문제가 없고, 결합법칙과 분배법칙은 다항식에서처럼 성립한다. 다항식은 계수가 유한개만 $0$이 아닌 형식적 멱급수이다. 우리가 쓰는 계수환은 정수 $\mathbb Z$나 유리수 $\mathbb Q$이다.

\begin{proposition}
\label{prop:fps-inverse}
$A=\sum_ka_kq^k\in K[[q]]$가 곱셈의 역원을 가질 필요충분조건은 $a_0$이 $K$의 가역원인 것이다. 특히 양의 정수 $j$에 대해 $\mathbb Z[[q]]$에서
\[
\frac1{1-q^j}=\sum_{m\ge0}q^{jm}=1+q^j+q^{2j}+\cdots
\]
이다.
\end{proposition}

\begin{proof}
$AB=1$이면 상수항에서 $a_0b_0=1$이므로 $a_0$은 가역원이다. 거꾸로 $a_0$이 가역원이면 $b_0=a_0^{-1}$로 두고 $s\ge1$에 대해 $b_s=-a_0^{-1}\sum_{k=1}^sa_kb_{s-k}$로 차례로 정의하자. 그러면 $[q^s](AB)=a_0b_s+\sum_{k=1}^sa_kb_{s-k}=0$이 모든 $s\ge1$에서 성립하므로 $AB=1$이다. 둘째 주장은 $(1-q^j)\sum_mq^{jm}$을 전개하면 $q^{jm}$들이 이웃끼리 지워져 $1$만 남는다는 것으로 확인된다.
\end{proof}

이 증명의 점화식은 그대로 알고리즘이다. $b_s$를 정하는 데는 $a_0,\dots,a_s$와 $b_0,\dots,b_{s-1}$만 필요하므로, 역원의 처음 $s+1$개 계수는 $A$의 처음 $s+1$개 계수만으로 정해진다. 같은 성질을 곱셈에 대해 정리해 두자.

\begin{proposition}
\label{prop:fps-truncate}
$\pi_b:K[[q]]\to K[[q]]/(q^{b+1})$를 차수 $b$ 이하의 계수만 남기는 사상이라 하면 $\pi_b$는 환 준동형사상이다. 따라서 곱셈과 가역원의 역원을 거듭해 얻는 어떤 식이든, 각 단계에서 차수 $b$ 위를 버리고 계산해도 마지막 결과의 차수 $b$ 이하 계수는 버리지 않고 계산한 것과 같다.
\end{proposition}

\begin{proof}
$[q^s](AB)$는 $a_0,\dots,a_s$와 $b_0,\dots,b_s$에만 의존하므로 $\pi_b(AB)$는 $\pi_b(A)$와 $\pi_b(B)$로 정해지고, 덧셈은 계수별이다. 따라서 $\pi_b$는 준동형이다. $\pi_b(A^{-1})=\pi_b(A)^{-1}$은 준동형이 단위원을 보존한다는 데서 나온다.
\end{proof}

검정의 p값에는 분포 전체가 필요하지 않고 관측값 쪽 꼬리만 필요하다. \cref{prop:fps-truncate}는 그 꼬리가 끝나는 차수 $b$에서 모든 다항식을 잘라 계산해도 된다는 보증이며, 이 장의 알고리즘들이 일의 양을 꼬리의 길이에 비례하게 줄이는 근거이다.

\subsection{인수 하나를 곱하고 나누는 알고리즘}

이 장에서 만나는 생성함수는 대부분 $(1-q^a)$꼴 인수들의 곱과 몫이므로, 계수 배열 $P[0..b]$에 이런 인수 하나를 곱하거나 나누는 연산이 기본 단위가 된다. $(1-q^a)$를 곱하면 $[q^v]$ 계수가 $P[v]-P[v-a]$로 바뀌는데, 제자리에서 계산하려면 $v$를 위에서부터 내려와야 $P[v-a]$가 아직 바뀌지 않은 값으로 남는다. \cref{prop:fps-inverse}에서 $(1-q^j)$로 나누는 것은 $\sum_mq^{jm}$을 곱하는 것이고, 결과 $Q$는 $Q(1-q^j)=P$, 곧 $Q[v]=P[v]+Q[v-j]$를 만족하므로 이번에는 $v$를 아래에서부터 올라가며 $P[v]\gets P[v]+P[v-j]$를 하면 $P[v-j]$가 이미 새 값 $Q[v-j]$로 바뀌어 있다. 나눗셈이 보폭 $j$의 누적합이 된다. 두 연산 모두 차수 $b$까지의 덧셈 또는 뺄셈 $b$번 남짓이다.

\begin{algorithm}[t]
\begin{algorithmic}[1]
\Require 계수 배열 $P[0..b]$, 양의 정수 $a$, $j$
\Ensure $P\gets P\cdot(1-q^a)/(1-q^j)\bmod q^{b+1}$
\For{$v\gets b$ \textbf{downto} $a$} \Comment{$(1-q^a)$를 곱한다}
  \State $P[v]\gets P[v]-P[v-a]$
\EndFor
\For{$v\gets j$ \textbf{to} $b$} \Comment{$(1-q^j)$로 나눈다: 보폭 $j$의 누적합}
  \State $P[v]\gets P[v]+P[v-j]$
\EndFor
\end{algorithmic}
\caption{$(1-q^a)/(1-q^j)$를 제자리에서 곱하기}
\label{alg:factor-step}
\end{algorithm}

\subsection{분할}

생성함수는 정수의 분할에서 처음으로 큰 힘을 보였다. 양의 정수 $s$를 순서를 무시하고 양의 정수들의 합으로 쓰는 방법을 $s$의 \term{분할}{partition}이라 하고 그 수를 $p(s)$로 쓴다. 예컨대 $4=3+1=2+2=2+1+1=1+1+1+1$이므로 $p(4)=5$이다. 분할은 각 크기 $k$를 몇 번 쓰는지 $m_k\ge0$을 정하는 것과 같고, 크기 $k$를 $m_k$번 쓰면 합에 $km_k$가 보태지므로
\[
\sum_{s\ge0}p(s)q^s=\prod_{k\ge1}\bigl(1+q^k+q^{2k}+\cdots\bigr)=\prod_{k\ge1}\frac1{1-q^k}
\]
이다. 곱의 각 인수가 ``크기 $k$를 몇 번 쓸지''라는 독립적인 선택 하나를 나타내고 곱셈이 그 선택들을 모두 조합한다. 쓸 수 있는 부분의 크기를 $n$ 이하로 제한하면 곱이 $k\le n$까지로 줄어들고, 부분의 개수를 제한하는 것도 같은 꼴로 다룰 수 있다. 다음 절의 $q$-이항계수는 부분의 개수와 크기를 동시에 제한한 분할을 센다.

\begin{aside}[Euler와 분할]
Euler는 1740년 무렵 Philippe Naudé가 보낸 편지에서 정수를 정해진 개수의 부분으로 나누는 방법의 수를 묻는 질문을 받고, 그 답을 무한곱의 전개로 구하는 방법을 고안했다. 그는 이 방법으로 서로 다른 부분으로 나눈 분할의 수와 홀수 부분으로 나눈 분할의 수가 같다는 것을 $\prod(1+q^k)=\prod(1-q^{2k})/(1-q^k)=\prod1/(1-q^{2k-1})$라는 한 줄로 보였는데, 같은 사실을 대응 관계로 직접 증명하는 것은 훨씬 까다롭다. 오늘날의 생성함수 방법은 수열을 멱급수에 실어 대수로 셈하는 이 발상에서 출발한다\cite{andrews1998partitions}.
\end{aside}

\section{가우스 이항계수}
\label{sec:qbinom}

\subsection{정의와 $q$-Pascal 공식}

정수 $m\ge0$에 대해 $[m]_q=1+q+\cdots+q^{m-1}=(1-q^m)/(1-q)$로 두고, $[m]_q!=[1]_q[2]_q\cdots[m]_q$($[0]_q!=1$)라 하자. $q\to1$이면 $[m]_q\to m$이므로 이것들은 정수와 계승의 $q$-판이다.

\begin{definition}[가우스 이항계수]
\label{def:qbinom}
$0\le k\le n$에 대해
\[
\qbinom nk=\frac{[n]_q!}{[k]_q!\,[n-k]_q!}=\frac{(1-q^n)(1-q^{n-1})\cdots(1-q^{n-k+1})}{(1-q)(1-q^2)\cdots(1-q^k)}
\]
를 \term{가우스 이항계수}{Gaussian binomial coefficient} 또는 $q$-이항계수라 하고, 그 밖의 $k$에서는 $0$으로 둔다.
\end{definition}

정의는 $\mathbb Q(q)$의 유리함수이지만 실제로는 음이 아닌 정수 계수의 다항식이다. 이 사실은 다음의 $q$-Pascal 공식에서 귀납법으로 바로 나온다.

\begin{proposition}
\label{prop:qpascal}
$1\le k\le n$에 대해
\begin{equation}
\label{eq:qpascal}
\qbinom nk=\qbinom{n-1}k+q^{n-k}\qbinom{n-1}{k-1}
\end{equation}
이다. 따라서 $\qbinom nk$는 음이 아닌 정수 계수를 갖는 차수 $k(n-k)$의 다항식이고 $q=1$에서 $\binom nk$와 같다.
\end{proposition}

\begin{proof}
우변의 두 항을 공통 인수 $[n-1]_q!/([k]_q!\,[n-k]_q!)$로 묶으면
\[
\qbinom{n-1}k+q^{n-k}\qbinom{n-1}{k-1}=\frac{[n-1]_q!}{[k]_q!\,[n-k]_q!}\bigl([n-k]_q+q^{n-k}[k]_q\bigr)
\]
이다. 괄호 안은 $(1+\cdots+q^{n-k-1})+(q^{n-k}+\cdots+q^{n-1})=[n]_q$이므로 우변은 $\qbinom nk$이다. $\qbinom n0=\qbinom nn=1$에서 시작해 $n$에 대한 귀납법을 쓰면 계수가 음이 아닌 정수라는 것과 차수가 $\max\bigl(k(n-1-k),\,(n-k)+(k-1)(n-k)\bigr)=k(n-k)$라는 것을 얻고, $q=1$을 넣으면 \cref{eq:pascal}로 돌아간다.
\end{proof}

괄호 안을 $[k]_q+q^k[n-k]_q=[n]_q$로 묶는 같은 계산으로 짝이 되는 공식 $\qbinom nk=q^k\qbinom{n-1}k+\qbinom{n-1}{k-1}$도 얻는다. 또 정의에서 $\qbinom nk=\qbinom n{n-k}$는 분명하다.

\subsection{조합적 의미}

$\{1,\dots,n\}$의 $k$원소 부분집합 $S=\{s_1<s_2<\cdots<s_k\}$에 대해
\[
a(S)=\sum_{i=1}^k(s_i-i)
\]
라 하자. $S$를 길이 $n$의 $0$과 $1$의 낱말 $w$로 바꾸어 $s_i$번째 글자만 $1$로 두면, $s_i-i$는 $i$번째 $1$보다 앞에 있는 $0$의 개수이므로 $a(S)$는 $0$이 $1$보다 앞에 오는 쌍 $(0,1)$의 수이다. 또 $s_i-i$들은 $0\le s_1-1\le s_2-2\le\cdots\le s_k-k\le n-k$를 만족하므로, $S$는 부분이 $k$개 이하이고 각 부분이 $n-k$ 이하인 분할(크기 $a(S)$)과 일대일로 대응한다. 그림으로는 $k\times(n-k)$ 상자 안에 들어가는 Young 그림이다.

\begin{theorem}
\label{thm:qbinom-count}
\[
\qbinom nk=\sum_{\substack{S\subseteq\{1,\dots,n\}\\|S|=k}}q^{a(S)}
\]
이다. 곧 $[q^u]\qbinom nk$는 $0$이 $1$보다 앞에 오는 쌍이 정확히 $u$개인, $1$이 $k$개 든 길이 $n$의 낱말의 수이다.
\end{theorem}

\begin{proof}
우변을 $G(n,k)$라 하자. $G(n,0)=G(n,n)=1$이고, $1\le k<n$일 때 $n\in S$인지로 경우를 나눈다. $n\notin S$이면 $S$는 $\{1,\dots,n-1\}$의 $k$원소 부분집합이고 $a(S)$는 그대로이므로 이런 $S$의 기여는 $G(n-1,k)$이다. $n\in S$이면 $s_k=n$이고 그 항이 $a(S)$에 $n-k$를 보태며, 나머지 $S\setminus\{n\}$은 $\{1,\dots,n-1\}$의 $(k-1)$원소 부분집합이므로 기여는 $q^{n-k}G(n-1,k-1)$이다. 따라서 $G$는 \cref{eq:qpascal}과 같은 점화식과 같은 경곗값을 가지므로 $G(n,k)=\qbinom nk$이다.
\end{proof}

이 정리에서 대칭성 $[q^u]\qbinom nk=[q^{k(n-k)-u}]\qbinom nk$도 곧바로 읽힌다. 낱말을 거꾸로 읽고 $0$과 $1$을 바꾸면 $1$의 개수는 그대로 $k$이고 $(0,1)$ 쌍과 $(1,0)$ 쌍의 역할이 바뀌는데, 두 종류의 쌍은 합쳐서 $k(n-k)$개이기 때문이다. 계수가 가운데까지 증가하다가 감소한다는 단봉성도 성립하지만(Sylvester가 1878년에 증명했다) 증명은 훨씬 어렵고 이 책에서는 쓰지 않는다.

\subsection{Mann--Whitney 통계량의 정확 분포}

크기 $c$인 표본 $x_1,\dots,x_c$와 크기 $N-c$인 표본 $y_1,\dots,y_{N-c}$를 합쳐 정렬하고, $x$들이 받은 순위를 $r_1<\cdots<r_c$라 하자. 동점이 없을 때
\[
U=\#\{(i,j):x_i>y_j\}=\sum_{i=1}^cr_i-\frac{c(c+1)}2
\]
를 \term{Mann--Whitney 통계량}{Mann--Whitney statistic}이라 한다\cite{mannwhitney1947}. 두 식이 같은 이유는 $r_i-i$가 $i$번째로 작은 $x$보다 작은 $y$의 개수이기 때문이고, 순위합에서 상수를 뺀 꼴이므로 Wilcoxon의 순위합 검정\cite{wilcoxon1945}과 같은 검정이다. 순위의 집합 $S=\{r_1,\dots,r_c\}$에 대해 $U=a(S)$이다.

두 표본이 같은 연속 분포에서 독립으로 나왔다는 귀무가설 아래에서는 $N$개 값의 순서를 이루는 $N!$가지 배치가 모두 같은 확률이고, 따라서 $x$들이 차지하는 자리의 집합 $S$는 $\binom Nc$가지가 모두 같은 확률이다. \cref{thm:qbinom-count}에서 곧바로 다음을 얻는다.

\begin{corollary}
\label{cor:mwu}
귀무가설 아래에서 $U$의 확률생성함수는
\[
E[q^U]=\qbinom Nc\Bigm/\binom Nc
\]
이다. 곧 $P(U=u)=[q^u]\qbinom Nc\big/\binom Nc$이다.
\end{corollary}

이 장을 연 예는 $N=7$, $c=3$, $S=\{1,2,4\}$였으므로 $U=0+0+1=1$이다. $\qbinom73$의 계수는 \cref{tab:qbinom73}의 첫 줄과 같고 합이 $\binom73=35$이므로 $P(U\le1)=(1+1)/35=2/35$로 처음의 계산이 확인된다. 분포가 $c(N-c)/2$를 중심으로 대칭이므로 위쪽 꼬리는 $P(U\ge u)=P(U\le c(N-c)-u)$로 아래쪽 꼬리로 바꿀 수 있고, 따라서 언제나 관측값에 가까운 쪽, 곧 길이가 $c(N-c)/2$ 이하인 아래쪽 꼬리 하나만 세면 된다.

Mann과 Whitney는 원래 논문에서 이 분포를 확률의 점화식으로 계산했다. 가장 큰 관측값이 $x$일 확률은 $c/N$이고, 그러면 그 $x$가 모든 $y$보다 크므로 $U$에 $N-c$가 보태진 뒤 나머지는 크기 $(c-1,N-c)$의 같은 문제이다. 가장 큰 값이 $y$이면 $U$에 보태지는 것이 없다. 크기 $(c,d)$에서 $U=u$일 확률을 $f(c,d,u)$로 쓰면
\[
f(c,d,u)=\frac c{c+d}\,f(c-1,d,u-d)+\frac d{c+d}\,f(c,d-1,u)
\]
이다\cite{mannwhitney1947}. 우변이 두 확률의 볼록결합이라 뺄셈이 없다는 점은 뒤에서 다시 중요해진다. 다만 이 점화식은 $(c,d,u)$의 세 첨자 전체에 걸친 표를 채워야 하므로, 한 표본이 작을 때에는 다음의 곱 공식이 훨씬 빠르다.

\subsection{곱 공식으로 계수 구하기}

\cref{def:qbinom}의 둘째 꼴에서 이웃한 두 계수 사이의 관계
\begin{equation}
\label{eq:qbinom-product}
\qbinom nj=\qbinom n{j-1}\cdot\frac{1-q^{n-j+1}}{1-q^j}
\end{equation}
를 얻는다. 따라서 $P=1$에서 출발해 $j=1,\dots,c$에 대해 \cref{alg:factor-step}을 $a=N-j+1$로 되풀이하면 $P=\qbinom Nc$가 된다. 꼬리의 끝 $b$에서 잘라도 된다는 것은 \cref{prop:fps-truncate}가 보증하고, 일의 양은 $c(b+1)$번 남짓의 정수 덧셈과 뺄셈이다. \cref{tab:qbinom73}은 $N=7$, $c=3$에서 이 과정을 따라간다. $(1-q^a)$를 곱한 직후의 배열에는 음수가 나타나지만, 이어지는 나눗셈 뒤에는 다시 $\qbinom7j$의 음이 아닌 계수가 된다.

\begin{table}[t]
\begin{sdtable}{colspec={lccccccccccccc}, row{1}={font=\small\bfseries}}
\toprule
단계 & 0 & 1 & 2 & 3 & 4 & 5 & 6 & 7 & 8 & 9 & 10 & 11 & 12 \\
\midrule
$\times(1-q^7)$ & 1 & 0 & 0 & 0 & 0 & 0 & 0 & $-1$ & 0 & 0 & 0 & 0 & 0 \\
$\div(1-q)$ & 1 & 1 & 1 & 1 & 1 & 1 & 1 & 0 & 0 & 0 & 0 & 0 & 0 \\
$\times(1-q^6)$ & 1 & 1 & 1 & 1 & 1 & 1 & 0 & $-1$ & $-1$ & $-1$ & $-1$ & $-1$ & $-1$ \\
$\div(1-q^2)$ & 1 & 1 & 2 & 2 & 3 & 3 & 3 & 2 & 2 & 1 & 1 & 0 & 0 \\
$\times(1-q^5)$ & 1 & 1 & 2 & 2 & 3 & 2 & 2 & 0 & 0 & $-2$ & $-2$ & $-3$ & $-2$ \\
$\div(1-q^3)$ & 1 & 1 & 2 & 3 & 4 & 4 & 5 & 4 & 4 & 3 & 2 & 1 & 1 \\
\bottomrule
\end{sdtable}
\caption{곱 공식 \cref{eq:qbinom-product}로 $\qbinom73$을 만드는 과정. 열의 머리는 $q$의 차수이고, 둘째·넷째·마지막 줄은 각각 $\qbinom71$, $\qbinom72$, $\qbinom73$의 계수이다. 마지막 줄의 합은 $\binom73=35$이다.}
\label{tab:qbinom73}
\end{table}

정확한 정수로 계산할 때는 중간값의 크기를 따져 두어야 한다. 나눗셈의 누적합을 아래에서부터 진행하므로 갱신이 끝난 칸은 이미 $\qbinom Nj$의 계수이고, 갱신 전의 칸은 $\qbinom N{j-1}(1-q^a)$의 계수로서 절댓값이 $\qbinom N{j-1}$의 가장 큰 계수 이하이다. 계수들은 모두 $\binom Nj$나 $\binom N{j-1}$ 이하이므로, $c$를 작은 쪽 표본의 크기로 잡으면 모든 중간값과 꼬리의 합이 $\binom Nc$ 이하이다. 부호 있는 $128$비트 정수는 $2^{127}$ 미만을 담으므로 $\binom Nc<2^{127}$이면 계산 전체가 정확하다. 이 범위는 꽤 넓어서, 한 표본이 $8$개일 때 $\binom N8<2^{124}$인 $N$은 $174447$까지이므로 다른 표본이 $17$만 개를 넘어도 정확한 정수로 센다.

\begin{instardust}[\code{csrc/htest/rank\_exact.hpp}]
\code{qbinom\_step}은 \cref{alg:factor-step}의 두 반복문 그대로이고, \code{mwu\_lower}는 $j=1,\dots,c$에 대해 이것을 부른 뒤 차수 $b$까지의 계수를 더해 $\#\{U\le b\}$를 낸다. p값은 이 개수와 $\binom Nc$의 몫 하나이다. 정수 길은 여유를 두어 경우의 수가 $2^{124}$ 아래일 때만 쓰고, 그 위나 동점이 있는 경우는 \cref{sec:ties}의 double-double 확률로 넘어간다.
\end{instardust}

\subsection{$q$-이항정리}

이항정리의 $q$-판은 다음과 같다.

\begin{theorem}[$q$-이항정리]
\label{thm:qbinomial-theorem}
\[
\prod_{i=0}^{n-1}\bigl(1+q^it\bigr)=\sum_{k=0}^nq^{k(k-1)/2}\qbinom nk\,t^k.
\]
\end{theorem}

\begin{proof}
좌변을 전개하면 각 인수에서 $1$과 $q^it$ 중 하나를 고르므로, $q^it$를 고른 첨자들의 집합 $T=\{t_1<\cdots<t_k\}\subseteq\{0,\dots,n-1\}$마다 항 $q^{t_1+\cdots+t_k}t^k$가 하나씩 나온다. $S=\{t_1+1,\dots,t_k+1\}$로 두면 $\sum_it_i=\sum_i(s_i-i)+\sum_i(i-1)=a(S)+k(k-1)/2$이므로, $t^k$의 계수는 $q^{k(k-1)/2}\sum_{|S|=k}q^{a(S)}$이고 \cref{thm:qbinom-count}에서 이것은 $q^{k(k-1)/2}\qbinom nk$이다.
\end{proof}

$t=q$를 넣으면 $\prod_{i=1}^n(1+q^i)=\sum_kq^{k(k+1)/2}\qbinom nk$를 얻는다. 좌변은 $\{1,\dots,n\}$의 부분집합을 원소의 합으로 센 생성함수이고, 우변은 그것을 원소의 개수 $k$로 나누어 센 것이다. 이 식은 다음 절의 부호 순위 통계량과 Mann--Whitney 통계량이 같은 뿌리에서 나왔음을 보여 준다.

같은 생각으로 다른 순위 점수도 다룰 수 있다. 척도의 차이를 보는 Ansari--Bradley 검정\cite{ansari1960}은 합친 표본에서 순위 $r$에 점수 $\min(r,N+1-r)$을 주고 한 표본의 점수합을 통계량으로 쓴다. $N=2k$이면 점수는 $1,\dots,k$가 두 번씩 나타나는데, 크기 $c$인 표본이 첫째 사본에서 $j$개, 둘째 사본에서 $c-j$개를 차지한다면, 한 사본 안에서 $j$개를 고른 점수합은 $j(j+1)/2$에 $\qbinom kj$를 따르는 양을 더한 것이다(\cref{thm:qbinom-count}). 따라서 점수합의 생성함수는
\[
\sum_jq^{j(j+1)/2+(c-j)(c-j+1)/2}\qbinom kj\qbinom k{c-j}
\]
이고, $N=2k+1$이면 가운데 점수 $k+1$을 고르는 경우가 $q^{k+1}$을 곱한 같은 꼴의 항으로 더해진다. 모든 $j$에 대해 $\qbinom kj$를 한 번씩만 만들어 두면 꼬리는 두 계수열의 합성곱을 누적합과 짝지어 얻는다.

\section{부분집합의 합}
\label{sec:signed-rank}

짝지은 관측 $(x_i,y_i)$의 차 $d_i=x_i-y_i$가 $0$이 아닌 $n$개라 하고, $|d_i|$들의 순위를 매겨 차가 양수인 것들의 순위만 더한 값
\[
T^+=\sum_{i:\,d_i>0}\operatorname{rank}|d_i|
\]
를 \term{Wilcoxon 부호 순위 통계량}{Wilcoxon signed-rank statistic}이라 한다\cite{wilcoxon1945}. 차의 분포가 $0$에 대해 대칭이라는 귀무가설 아래에서는 각 차의 부호가 크기와 독립이고 반반의 확률로 양과 음이므로, 순위 $1,\dots,n$이 각각 독립적으로 반의 확률로 합에 들어간다. 따라서
\begin{equation}
\label{eq:signed-rank-gf}
E[q^{T^+}]=\prod_{i=1}^n\frac{1+q^i}2
\end{equation}
이고, $2^nP(T^+=s)$는 합이 $s$인 $\{1,\dots,n\}$의 부분집합의 수이다. $n=5$이면 $s=0,1,\dots,15$에 대해 이 수는
\[
1,1,1,2,2,3,3,3,3,3,3,2,2,1,1,1
\]
이고 합은 $2^5=32$이다. $n=10$에서 $P(T^+\le8)=25/1024\approx0.0244$인데, 이것은 열 쌍 가운데 차가 양수인 쌍의 순위합이 $8$ 이하일 확률이다.

\cref{eq:signed-rank-gf}의 전개는 전산학에서 \term{부분집합 합}{subset sum} 문제라 부르는 것의 세기 판이고 동적 계획으로 푼다. 계수 배열을 $C[0]=1$, 나머지 $0$으로 두고 인수 $(1+q^a)$를 하나씩 곱하면 되는데, 곱셈은 $C[s]\gets C[s]+C[s-a]$이다. 이것을 제자리에서 하려면 $s$를 위에서부터 내려와야 하는데, 아래에서부터 올라가면 이미 $q^a$가 곱해진 값을 다시 써서 원소 $a$를 두 번 넣는 셈이 되기 때문이다. 이 내려가는 순서는 각 원소를 한 번만 쓰는 0-1 배낭 문제의 표준 기법과 같다.

\begin{algorithm}[t]
\begin{algorithmic}[1]
\Require 양의 정수 점수 $a_1,\dots,a_n$
\Ensure $C[s]=\#\{I\subseteq\{1,\dots,n\}:\sum_{i\in I}a_i=s\}$, $s=0,\dots,\sum_ia_i$
\State $C[0]\gets1$, 나머지 $C[s]\gets0$; $\mathit{top}\gets0$
\For{$i\gets1$ \textbf{to} $n$}
  \State $\mathit{top}\gets\mathit{top}+a_i$
  \For{$s\gets\mathit{top}$ \textbf{downto} $a_i$} \Comment{내려가야 $a_i$를 한 번만 쓴다}
    \State $C[s]\gets C[s]+C[s-a_i]$
  \EndFor
\EndFor
\end{algorithmic}
\caption{부분집합 합의 개수}
\label{alg:subset-sum}
\end{algorithm}

점수가 $1,\dots,n$이면 합의 범위가 $n(n+1)/2$이므로 일은 $O(n^3)$이다. 경우의 수 $2^n$이 정수형을 넘는 큰 $n$에서는 개수 대신 확률을 쌓는다. 곱하는 인수가 $(1+q^a)/2$이므로 갱신은 $x(s)\gets\bigl(x(s)+x(s-a)\bigr)/2$가 되는데, 이진 부동소수점에서 $2$로 나누는 것은 아래로 넘치지 않는 한 정확하고, 두 양수의 합은 상대오차 하나만 보탠다. 그래서 $n$단계 뒤의 상대오차는 단위 반올림의 $n$배 정도로 묶인다.

대칭을 쓰면 일을 반으로 줄일 수 있다. 모든 부호를 뒤집으면 $T^+$가 $\sum a_i-T^+$로 바뀌므로 분포는 $\sum a_i/2$에 대해 대칭이고, 이것은 인수를 몇 개만 곱한 중간 단계에서도 그 단계의 점수합의 반에 대해 성립한다. 따라서 각 단계에서 배열의 아래쪽 반만 들고 있어도 된다. 갱신에 필요한 값 $x(s)$가 들고 있는 범위의 위에 있으면, 직전 단계의 점수합을 $\sigma$라 할 때 대칭에서 $x(s)=x(\sigma-s)$이고 $\sigma-s$는 들고 있는 범위 안에 있기 때문이다. 게다가 p값에는 관측값까지의 꼬리만 필요하므로 반쪽의 끝과 꼬리의 끝 가운데 작은 쪽까지만 계산한다.

\begin{instardust}[\code{csrc/htest/rank\_exact.hpp}]
\code{signed\_rank\_counts}는 \cref{alg:subset-sum}을 $128$비트 정수로 한 것이고, \code{signed\_rank\_lower}는 확률 판이다. 확률 판은 값마다 double-double을 들고 $x(s)\gets(x(s)+x(s-a))/2$를 아래쪽 반에서만 계산하며, 위의 범위에 필요한 값은 대칭인 자리에서 읽는다. 한 번에 네 칸씩 NEON 레지스터로 처리하고, 범위가 아주 길면 같은 연산을 조각으로 나누어 스레드풀에서 돌리되 다른 배열에 쓰게 해서 스레드 수와 상관없이 같은 비트를 낸다. 동점이 있으면 점수가 두 배 한 중간 순위라는 정수가 될 뿐 알고리즘은 같다.
\end{instardust}

\section{역순의 수}
\label{sec:inversions}

\subsection{순열의 역순과 Mahonian 수}

$\{1,\dots,n\}$의 순열 $\pi=(\pi_1,\dots,\pi_n)$에서 $i<j$이면서 $\pi_i>\pi_j$인 쌍을 \term{역순}{inversion}이라 하고 그 수를 $\operatorname{inv}\pi$로 쓴다. 쌍을 이룬 관측 $(x_i,y_i)$를 $x$의 순서로 정렬했을 때 $y$의 순위가 이루는 순열의 역순은 $x$와 $y$의 순서가 엇갈리는 \textbf{불일치 쌍}의 수 $K$이고, 동점이 없으면 Kendall의 순위상관계수는 $\tau=1-4K/(n(n-1))$이다\cite{kendall1938}. $x$와 $y$가 독립인 연속 확률변수라는 귀무가설 아래에서는 $n!$개의 순열이 모두 같은 확률이므로, $\tau$의 분포를 구하는 것은 역순의 수의 분포를 구하는 것이다.

\begin{theorem}
\label{thm:mahonian}
\[
\sum_{\pi\in S_n}q^{\operatorname{inv}\pi}=\prod_{j=1}^n[j]_q=[n]_q!.
\]
\end{theorem}

\begin{proof}
순열 $\pi$에 대해 $c_j=\#\{i<j:\pi_i>\pi_j\}$라 하면 $0\le c_j\le j-1$이고 $\operatorname{inv}\pi=\sum_jc_j$이다. 거꾸로 $0\le c_j\le j-1$인 수열 $(c_1,\dots,c_n)$이 주어지면 $\pi$를 오른쪽 끝부터 하나로 정할 수 있다. 위치 $n$의 값 $\pi_n$보다 왼쪽에 있으면서 큰 값이 $c_n$개이고 $\pi_n$보다 큰 값은 모두 $\pi_n$의 왼쪽에 있으므로, $\pi_n$은 전체에서 $(c_n+1)$번째로 큰 값이다. 그 값을 지우면 남은 $n-1$개의 값과 위치에 같은 논리가 적용되어, 위치 $j$의 값은 남은 $j$개 가운데 $(c_j+1)$번째로 큰 값이다. $c_j\le j-1$이므로 이 선택은 언제나 가능하다. 따라서 $\pi\mapsto(c_1,\dots,c_n)$은 $S_n$에서 $\prod_j\{0,\dots,j-1\}$로 가는 전단사이고, 생성함수는 $\prod_j(1+q+\cdots+q^{j-1})$이다.
\end{proof}

증명의 수열 $(c_j)$를 \term{Lehmer 부호}{Lehmer code}라 한다. 계수 $M_n(k)=[q^k][n]_q!$를 \term{Mahonian 수}{Mahonian numbers}라 하는데 MacMahon의 이름을 딴 것으로, 그는 역순의 수와 major index라는 다른 통계량이 같은 분포를 가진다는 것을 보였다\cite{macmahon1915}. 생성함수 자체는 그보다 훨씬 앞서 Rodrigues가 1839년에 얻었다\cite{knuth1998taocp3}. $n=4$와 $n=5$의 Mahonian 수는
\[
1,3,5,6,5,3,1\qquad\text{와}\qquad1,4,9,15,20,22,20,15,9,4,1
\]
이고 합은 각각 $4!=24$와 $5!=120$이다.

\subsection{창의 합}

\cref{thm:mahonian}에서 $[j]_q=(1-q^j)/(1-q)$를 하나씩 곱하는 것이 점화식
\[
M_j(k)=\sum_{i=0}^{j-1}M_{j-1}(k-i)
\]
이다. 우변은 길이 $j$의 \term{미끄러지는 창}{sliding window} 위의 합이고, 정수라면 $k$를 하나씩 늘리며 들어오는 항을 더하고 나가는 항을 빼는 누적값 하나로 단계마다 $O(1)$에 계산한다. 이것은 \cref{alg:factor-step}에서 $(1-q^j)$를 곱한 뒤 $(1-q)$로 나누는 것과 같은 계산이다. 꼬리 $k\le b$만 필요하므로 일은 $n(b+1)$이다. $n!<2^{127}$인 동안은 부호 있는 $128$비트 정수로 정확히 셀 수 있는데, $33!/2^{127}\approx0.051$이고 $34!/2^{127}\approx1.74$이므로 그 경계는 $n=33$이다.

더 큰 $n$에서는 확률 $g_j(k)=M_j(k)/j!$를 직접 쌓는다. 위 점화식을 $j!$로 나누면
\[
g_j(k)=\frac1j\sum_{i=0}^{j-1}g_{j-1}(k-i)
\]
인데, 이것은 $j$번째 원소가 들어갈 자리가 $j$곳 가운데 균등하게 정해지면서 역순을 $0$개에서 $j-1$개까지 보탠다는 확률적 해석 그대로이다. 그런데 부동소수점에서 미끄러지는 창의 누적값을 쓰면 문제가 생긴다. 누적값은 사실상 두 누적합의 차 $\sum_{i\le k}g-\sum_{i\le k-j}g$이고, 반올림 오차는 창의 합이 아니라 누적합의 크기에 비례해 쌓이므로, 창의 합이 누적합보다 훨씬 작은 자리에서는 상대 정확도를 잃는다. 뺄셈에서 생기는 이런 소거는 9장에서 자세히 다룬다.

뺄셈 없이 창의 합을 구하는 방법이 있다. 첨자를 길이 $j$의 블록 $[0,j)$, $[j,2j)$, $\dots$로 나누고, 블록마다 앞에서부터의 누적합 $\mathit{pre}$와 뒤에서부터의 누적합 $\mathit{suf}$를 만들어 두자. 길이가 $j$인 창 $[k-j+1,k]$는 블록 경계와 일치하거나 정확히 이웃한 두 블록에 걸친다. 앞의 경우 창의 합은 $\mathit{pre}[k]$이고, 뒤의 경우 창의 왼쪽 끝이 든 블록의 뒷부분과 오른쪽 끝이 든 블록의 앞부분의 합이므로 $\mathit{suf}[k-j+1]+\mathit{pre}[k]$이다. 블록 누적합을 만드는 데 칸마다 덧셈 두 번, 창 하나에 덧셈 한 번이면 되고 모든 덧셈은 양수끼리이다. 이 기법은 영상 처리에서 쓰는 van Herk와 Gil--Werman의 알고리즘과 같은 구조인데\cite{vanherk1992,gil1993}, 창 위의 최댓값과 최솟값을 원소마다 상수 번의 비교로 구하는 그 알고리즘에서는 최댓값에 역연산이 없어서 같은 고안이 필요했다.

\begin{algorithm}[t]
\begin{algorithmic}[1]
\Require 확률 배열 $g[0..b]$ ($g_{j-1}$), 창의 길이 $j$
\Ensure $g[k]\gets(1/j)\sum_{i=0}^{j-1}g[k-i]$, $k=0,\dots,b$ ($g_j$, 음수 첨자는 $0$)
\For{블록의 시작 $s\gets0,j,2j,\dots\le b$} \Comment{블록 $[s,\min(b,s+j-1)]$}
  \State $\mathit{pre}[s..e]\gets$ 앞에서부터의 누적합, $\mathit{suf}[s..e]\gets$ 뒤에서부터의 누적합
\EndFor
\For{$k\gets0$ \textbf{to} $b$}
  \State $l\gets k-j+1$ \Comment{창 $[l,k]$}
  \If{$l>0$ \textbf{and} $l\bmod j\ne0$}
    \State $w\gets\mathit{suf}[l]+\mathit{pre}[k]$ \Comment{이웃한 두 블록에 걸친 창}
  \Else
    \State $w\gets\mathit{pre}[k]$
  \EndIf
  \State $g[k]\gets w/j$
\EndFor
\end{algorithmic}
\caption{뺄셈 없는 창의 합으로 Mahonian 확률을 한 단계 넘기기}
\label{alg:block-window}
\end{algorithm}

확률을 쌓다 보면 마지막으로 아래로 넘침을 만난다. $P(K=0)=1/n!$이고 $n=200$이면 $1/200!\approx1.27\times10^{-375}$인데, 이것은 \code{double}의 가장 작은 서브노멀 수 약 $4.9\times10^{-324}$보다도 작다. 그래서 가장 큰 항이 $2^{-600}$ 아래로 내려가면 배열 전체에 $2^{600}$을 곱해 두고 곱한 횟수를 기억했다가 마지막에 되돌린다. $2$의 거듭제곱을 곱하는 것은 지수만 바꾸는 정확한 연산이므로 이 조정은 오차를 보태지 않으며, 마지막 결과는 \code{double}로 나타낼 수 있는 범위에서 올바르게 서브노멀이나 $0$이 된다.

\begin{instardust}[\code{csrc/htest/kendall.hpp}]
Kendall의 정확 분포는 $n\le33$이면 $128$비트 정수의 미끄러지는 창으로, 그 위는 \cref{alg:block-window}를 double-double로 한 것으로 계산한다. 주석은 설계 초안의 미끄러지는 창이 뺄셈 때문에 꼬리의 자릿수를 잃어서 블록 꼴로 바꾸었다고 밝혀 둔다. $j$로 나누는 것은 $j$가 $2$의 거듭제곱이면 정확하므로 그때는 두 부분에 $1/j$를 곱하기만 하고, 그 밖에는 double-double 나눗셈을 쓴다. 분포가 대칭이므로 관측값에 가까운 쪽 꼬리만 계산하며, 일의 양 $n(c+1)$이 정해진 상한을 넘으면 정확 계산 대신 정규 근사를 쓴다.
\end{instardust}

\section{동점이 있을 때}
\label{sec:ties}

\subsection{두 변수 생성함수}

관측값에 같은 값이 있으면 순위가 하나로 정해지지 않는다. 표준적인 처리에서는 같은 값들에 그들이 차지하는 순위의 평균인 \term{중간 순위}{midrank}를 준다. 위치 $b+1$부터 $b+t$까지를 같은 값 $t$개가 차지하면 중간 순위는 $b+(t+1)/2$이고, 그 두 배 $2b+t+1$은 정수이다. 이제 귀무분포는 관측한 점수들을 고정해 둔 \textbf{조건부} 분포로 정의한다. $N$개의 점수가 정해진 상태에서 한 표본이 그 가운데 어느 $c$개를 차지할지가 $\binom Nc$가지 모두 같은 확률이라고 본다.

점수가 서로 다른 값 $v_1,\dots,v_G$로 묶이고 값 $v_g$가 $t_g$번 나타난다고 하자. 묶음 $g$에서 $k$개를 고르는 방법은 $\binom{t_g}k$가지이고 점수합에 $kv_g$를 보태므로, 고른 개수를 $z$로, 점수합을 $q$로 표시한 두 변수 생성함수는
\[
F(z,q)=\prod_{g=1}^G\sum_{k=0}^{t_g}\binom{t_g}kz^kq^{kv_g}
\]
이고 $[z^cq^s]F$가 점수합이 $s$인 선택의 수이다. 동점이 없으면 모든 $t_g=1$이고 $F=\prod_g(1+zq^{v_g})$인데, 점수가 $0,1,\dots,N-1$이면 이것이 바로 \cref{thm:qbinomial-theorem}의 좌변이다.

전개는 고른 개수 $j$와 합 $s$를 첨자로 갖는 두 차원 배열 $D[j][s]$ 위의 동적 계획이다. 묶음을 하나씩 들이며, $j$를 위에서부터 내려오면서 그 묶음에서 $k\ge1$개를 더 고르는 경우를 $D[j][s+kv_g]\mathrel{+}=\binom{t_g}kD[j-k][s]$로 보탠다. 이번에도 내려가는 순서가 한 묶음을 두 번 쓰지 않게 한다. 일의 양은 대략 (전체 관측 수)$\times c\times$(합의 범위)이다. 정확히 $c$개를 고르므로 모든 점수에서 가장 작은 점수를 빼면 합이 $c\cdot v_{\min}$만큼 옮겨질 뿐이고, 점수들의 최대공약수로 나누면 합의 범위가 그만큼 줄어든다.

동점이 있으면 분포가 더는 대칭이 아닐 수 있다. 값 $1,1,2,3,4$의 두 배 중간 순위는 $3,3,6,8,10$인데, 이 가운데 두 개를 고르는 $10$가지의 점수합은
\[
6,\ 9,\ 9,\ 11,\ 11,\ 13,\ 13,\ 14,\ 16,\ 18
\]
이다. 평균은 $12$인데 $P(S\le9)=3/10$이고 $P(S\ge15)=2/10$이므로, 평균에서 같은 거리만큼 떨어진 두 꼬리의 확률이 다르다. 그래서 동점이 있으면 대칭을 빌려 한쪽 꼬리로 다른 쪽을 얻을 수 없고 두 꼬리를 따로 세어야 한다.

\begin{instardust}[\code{csrc/htest/rank\_exact.hpp}]
\code{grouped\_counts}가 위의 동적 계획이며, 묶음에서 $k$개를 고르는 수 $\binom{t_g}k$와 개수들을 모두 double-double로 든다. 모든 갱신이 양수의 덧셈이므로 단계마다의 상대오차는 약 $2^{-104}$이다. 한 행에서 다음 행으로 옮겨 더하는 \code{dd\_axpy}는 \code{vld2q\_f64}로 double-double 두 개의 높은 부분과 낮은 부분을 갈라 읽어 NEON 레지스터 하나에 두 칸씩 처리한다. 동점을 이렇게 조건부로 다루는 것은 scipy와 다른 선택이다. scipy의 \code{mannwhitneyu}는 동점이 있을 때 정확한 방법을 고르면 동점이 없는 분포에 통계량을 넣어 버리므로, 그 p값은 그 데이터의 귀무분포에서 나온 값이 아니다.
\end{instardust}

\subsection{행의 순열을 부분집합으로 세기}

Page의 순서 대립 검정\cite{page1963}은 $m$개의 블록(행) 각각에서 $n$개의 처리(열)에 순위를 매기고, 처리에 미리 정한 순서 $1,\dots,n$을 가중치로 곱한 $L=\sum_jj\,R_j$($R_j$는 열 $j$의 순위합)를 통계량으로 쓴다. 귀무가설 아래에서 각 행의 값들이 열에 놓이는 순서는 $n!$가지가 모두 같은 확률이고 행끼리는 독립이므로, 행 하나의 통계량 $\sum_jj\,r_{\pi(j)}$의 생성함수를 구해 $m$개를 곱하면 된다.

행 하나의 분포를 $n!$개의 순열을 나열해 구하면 $n=12$에서 이미 $4.8$억 가지이다. 그러나 순열을 왼쪽 열부터 채워 간다고 보면, 다음 열에 놓을 값을 고를 때는 지금까지 어떤 값들을 썼는지만 중요하고 그것들을 어떤 순서로 놓았는지는 상관없다. 이미 쓴 값의 집합 $A$가 정해지면 다음 열의 번호는 $|A|+1$이기 때문이다. 그래서 상태를 값들의 부분집합 $A$로 잡고 각 상태에 지금까지의 합의 분포를 붙인 배열 $D[A][s]$를 두면
\[
D[A\cup\{i\}][s+(|A|+1)r_i]\mathrel{+}=D[A][s]\qquad(i\notin A)
\]
로 상태를 $|A|$의 순서대로 채워 갈 수 있다. 부분집합은 $2^n$개이고 각각에서 다음 값을 $n$가지로 고르므로 일은 $2^n\cdot n\cdot$(합의 범위)이며, $n!$이 $2^n$으로 바뀐 것이 이 방법의 핵심이다. 부분집합을 $n$비트 정수의 비트로 나타내는 이런 동적 계획을 \term{비트마스크 동적 계획}{bitmask dynamic programming}이라 하고 순열 위의 최적화나 세기 문제에서 널리 쓰인다.

\begin{instardust}[\code{csrc/htest/rank\_exact.hpp}]
\code{page\_row\_counts}는 이 동적 계획을 $64$비트 정수로 하며, 같은 크기의 부분집합을 한 층으로 처리하고 다음 층으로 넘긴 상태의 배열은 바로 비워서 기억 공간을 아낀다. 개수의 합이 $n!$이므로 $n!<2^{63}$인 $n\le20$까지 정확하다. 행들의 분포는 double-double로 합성곱해 생성함수의 곱을 만들고, 동점이 없으면 모든 행의 분포가 같으므로 꼬리의 누적을 $(m,n)$마다 한 번 만들어 둔다.
\end{instardust}

\section{격자 경로와 반사 원리}
\label{sec:lattice}

\subsection{두 표본 Kolmogorov--Smirnov 통계량}

크기 $m$과 $n$인 두 표본의 경험 분포 함수를 $F_m$, $G_n$이라 할 때 $D=\sup_x|F_m(x)-G_n(x)|$를 두 표본 Kolmogorov--Smirnov 통계량이라 한다. 합친 표본을 정렬해 왼쪽부터 읽으면서 첫째 표본의 값을 만나면 오른쪽으로, 둘째 표본의 값을 만나면 위로 한 칸씩 가는 \term{격자 경로}{lattice path}를 그리면, 첫째 표본 $i$개와 둘째 표본 $j$개를 지난 점 $(i,j)$에서 두 함수의 차는 $i/m-j/n$이다. $L=\operatorname{lcm}(m,n)$, $\alpha=L/m$, $\beta=L/n$으로 두면 이 차는 $(i\alpha-j\beta)/L$로서 $1/L$의 정수배이고, 따라서 $D\ge h/L$인지는 경로가 정수 조건 $|i\alpha-j\beta|\ge h$를 만족하는 점을 지나는지로 판정된다. 실수의 비교를 정수의 비교로 바꾸었으므로 수학적으로 같은 두 값이 반올림 때문에 갈려 경우 하나를 빠뜨리거나 더 세는 일이 없다.

동점이 없고 귀무가설이 참이면 $\binom{m+n}m$개의 경로가 모두 같은 확률이므로, p값은 경계에 닿는 경로의 수를 전체 경로의 수로 나눈 값이다. 크기가 같은 $m=n$이면 $L=n$, $\alpha=\beta=1$이므로 경로를 높이 $g=i-j$가 한 걸음마다 $\pm1$씩 바뀌는 길이 $2n$의 걸음으로 볼 수 있고, 이 경우에는 개수를 닫힌 꼴로 얻는다.

\subsection{반사 원리}

\begin{figure}[t]
\centering
\begin{tikzpicture}[x=8mm, y=5mm, every node/.style={font=\small}]
  \draw[black!25, very thin] (0,-1) grid (10,6);
  \draw[->] (0,0) -- (10.5,0) node[right] {걸음};
  \draw[->] (0,-1) -- (0,6.5) node[above] {$g$};
  \draw[dashed] (0,3) -- (10,3) node[right] {$h$};
  \foreach \y in {0,3,6} \node[left] at (0,\y) {$\y$};
  \draw[thick] (0,0)--(1,1)--(2,2)--(3,1)--(4,2)--(5,3)--(6,2)--(7,1)--(8,0)--(9,1)--(10,0);
  \draw[thick, dashed] (5,3)--(6,4)--(7,5)--(8,6)--(9,5)--(10,6);
  \fill (5,3) circle (1.6pt);
  \fill (10,0) circle (1.6pt);
  \fill (10,6) circle (1.6pt);
\end{tikzpicture}
\caption{반사 원리. 실선은 $0$에서 출발해 $0$으로 돌아오는 길이 $10$의 경로로, 다섯째 걸음에서 처음으로 높이 $h=3$에 닿는다. 그 뒤의 부분을 직선 $g=h$에 대해 뒤집으면(점선) 끝점이 $2h=6$인 경로가 되고, 이 대응은 일대일이다.}
\label{fig:reflection}
\end{figure}

\begin{theorem}[반사 원리]
\label{thm:reflection}
양의 정수 $h\le n$에 대해, $\pm1$ 걸음 $2n$개로 $0$에서 $0$으로 가는 경로 가운데 높이 $h$에 닿는 것의 수는 $\binom{2n}{n-h}$이다.
\end{theorem}

\begin{proof}
높이 $h$에 닿는 경로가 처음으로 $h$에 닿는 걸음을 $\tau$라 하고, $\tau$ 뒤의 모든 걸음의 부호를 바꾸자. 바뀐 경로는 $\tau$ 뒤에 직선 $g=h$에 대해 원래 경로와 대칭이므로 끝점이 $h+(h-0)=2h$이다(\cref{fig:reflection}). 거꾸로 $0$에서 $2h$로 가는 경로는 $h>0$이므로 반드시 $h$를 지나고, 처음으로 $h$에 닿은 뒤의 걸음을 다시 뒤집으면 $0$으로 끝나면서 $h$에 닿는 경로가 된다. 두 변환은 서로의 역이므로 구하는 수는 $0$에서 $2h$로 가는 경로의 수이다. 그런 경로는 올라가는 걸음이 $n+h$개, 내려가는 걸음이 $n-h$개이므로 그 수는 $\binom{2n}{n+h}=\binom{2n}{n-h}$이다.
\end{proof}

이 정리를 확률로 바꾸면 $m=n$이고 동점이 없을 때 한쪽 통계량 $D^+=\sup(F_n-G_n)$에 대해
\[
P\bigl(D^+\ge h/n\bigr)=\binom{2n}{n-h}\Bigm/\binom{2n}n
\]
이다. 양쪽 통계량은 $h$와 $-h$ 둘 중 하나에 닿는 경로를 세야 하는데, 두 사건이 겹치므로 포함과 배제가 필요하다.

\begin{proposition}
\label{prop:two-sided-reflection}
같은 조건에서 높이 $h$나 $-h$ 가운데 적어도 하나에 닿는 경로의 수는
\[
2\sum_{k\ge1}(-1)^{k-1}\binom{2n}{n-kh}
\]
이다.
\end{proposition}

\begin{proof}
$A_k^+$를 시각의 순서대로 $h,-h,h,\dots$의 높이를 번갈아 $k$번 거치는(사이에 다른 방문이 있어도 된다) 경로의 집합이라 하고, 부호를 바꾼 $A_k^-$도 같은 식으로 정의하자. 먼저 $|A_k^\pm|=\binom{2n}{n-kh}$임을 보인다. $A_k^+$의 경로를 처음으로 $h$에 닿은 뒤로 직선 $g=h$에 대해 뒤집으면, 원래의 $-h$ 방문은 $3h$ 방문이 된다. 다시 처음으로 $3h$에 닿은 뒤로 직선 $g=3h$에 대해 뒤집는 식으로 $k$번 반사하면 끝점이 $2kh$인 경로를 얻는다. 끝점이 $2kh$인 경로는 반드시 $h,3h,\dots,(2k-1)h$를 차례로 지나므로 같은 반사를 거꾸로 적용할 수 있고, 따라서 이 대응은 전단사이다. 끝점이 $2kh$인 경로의 수는 $\binom{2n}{n-kh}$이다.

이제 경로마다 경계에 닿는 순서를 기록하되 같은 경계에 연이어 닿은 것은 하나로 합친 교대 수열을 생각하고, 그 길이를 $\ell$이라 하자. 이 수열이 $+h$로 시작하는 경로는 $k\le\ell$인 $A_k^+$와 $k\le\ell-1$인 $A_k^-$에 속하므로 $\sum_k(-1)^{k-1}(|A_k^+|+|A_k^-|)$에서 $\sum_{k=1}^\ell(-1)^{k-1}+\sum_{k=1}^{\ell-1}(-1)^{k-1}$번 세어진다. $\ell$이 홀수이면 첫 합이 $1$, 둘째 합이 $0$이고, 짝수이면 그 반대이므로 언제나 정확히 한 번이다. $-h$로 시작하는 경로도 마찬가지이고, 경계에 닿지 않는 경로는 어느 집합에도 속하지 않는다. 따라서 구하는 수는 $\sum_k(-1)^{k-1}\cdot2\binom{2n}{n-kh}$이다.
\end{proof}

Hodges는 두 표본 검정의 유의확률을 정리하면서 이 두 식을 닫힌 꼴로 썼다\cite{hodges1958}. $n=5$에서 직접 나열해 확인하면, $h=3$일 때 한쪽 확률은 $\binom{10}2/\binom{10}5=45/252=5/28\approx0.179$이고 양쪽 확률은 $2(45-0)/252=5/14\approx0.357$이다. $h=2$이면 양쪽 확률이 $2\bigl(\binom{10}3-\binom{10}1\bigr)/252=55/63$인데, 이 경우에는 두 경계에 모두 닿는 경로가 있어서 포함과 배제의 둘째 항이 실제로 필요하다. 계산할 때는 이항계수 대신 비
\[
\binom{2n}{n-l}\Bigm/\binom{2n}n=\prod_{i=0}^{l-1}\frac{n-i}{n+i+1}
\]
를 $l=h,2h,\dots$까지 차례로 곱해 가면 되므로 넘침이 없고 일은 $O(n)$이다. 교대합의 항들은 대략 $e^{-(kh)^2/n}$처럼 빠르게 줄어들므로 첫 항이 합을 지배하고, 뺄셈이 큰 소거를 일으키지 않는다.

\begin{aside}[투표 문제]
반사 원리는 1887년의 \term{투표 문제}{ballot problem}에서 나왔다. Bertrand는 후보 $A$가 $a$표, $B$가 $b$표($a>b$)를 얻었을 때 표를 무작위 순서로 하나씩 열면서 $A$가 처음부터 끝까지 앞설 확률이 $(a-b)/(a+b)$라는 결과를 냈고, 같은 해 André가 여기에 나쁜 순서를 뒤집어 세는 증명을 붙였다. 위 증명의 ``처음 닿은 뒤를 뒤집는다''는 생각이 이 전통에서 왔으며, Feller의 확률론 교과서가 이것을 무작위 걸음 이론의 출발점으로 삼은 뒤로 널리 퍼졌다\cite{feller1968}.
\end{aside}

\subsection{흡수 경계가 있는 동적 계획}

크기가 다르거나 동점이 있으면 닫힌 꼴이 없으므로 경로를 동적 계획으로 센다. 경계 안에 머무는 경로의 수를 세어 전체에서 빼는 방법이 가장 자연스럽지만 두 가지 이유로 좋지 않다. 경로의 수 $\binom{m+n}m$은 금세 \code{double}을 넘고, p값이 아주 작을 때 $1-P(\text{안})$을 계산하면 $1$에 가까운 수에서 $1$에 가까운 수를 빼는 소거로 자릿수를 거의 모두 잃는다. 이 때문에 경우의 수 대신 확률을 들고 경계 밖으로 나간 확률을 그 자리에서 따로 모은다.

작은 쪽 표본을 첫째 표본이라 하고 크기를 $m$이라 하자. 합친 표본에서 $p$개를 읽은 뒤의 상태를 그때까지 읽은 첫째 표본의 개수 $c$로 잡고, 상태 $c$에 ``지금까지 경계에 닿지 않고 이 상태에 올 확률'' $w(c)$를 붙인다. 남은 자리가 $R=N-p$개이고 그 가운데 첫째 표본이 $m-c$개이므로, 다음 값이 첫째 표본일 확률은 $(m-c)/R$이고 둘째 표본일 확률은 $(R-m+c)/R$이다. 한 걸음 옮긴 새 상태가 경계에 닿거나 넘으면 그 확률을 꼬리의 합 $\mathit{tail}$에 더하고 상태에서 지운다. 이렇게 하면 마지막 $\mathit{tail}$이 곧 $P(D\ge h/L)$이며, 모든 무게와 모든 덧셈이 양수이므로 상대오차가 걸음의 수에 비례하는 만큼만 쌓인다.

\begin{algorithm}[t]
\begin{algorithmic}[1]
\Require 크기 $m\le n$, $N=m+n$, 경계 $h$, 상태 $(c,p)$가 경계 밖인지 판정하는 함수 $\mathrm{out}(c,p)$
\Ensure $\mathit{tail}=P(D\ge h/L)$
\State $w[0]\gets1$, 나머지 $w[c]\gets0$; $\mathit{tail}\gets0$
\For{$p\gets0$ \textbf{to} $N-1$}
  \State $R\gets N-p$; $w'[\cdot]\gets0$
  \For{$w[c]>0$인 모든 $c$}
    \State $u\gets w[c]\cdot(m-c)/R$ \Comment{다음 값이 첫째 표본}
    \State \textbf{if} $\mathrm{out}(c+1,p+1)$ \textbf{then} $\mathit{tail}\gets\mathit{tail}+u$ \textbf{else} $w'[c+1]\gets w'[c+1]+u$
    \State $v\gets w[c]\cdot(R-m+c)/R$ \Comment{다음 값이 둘째 표본}
    \State \textbf{if} $\mathrm{out}(c,p+1)$ \textbf{then} $\mathit{tail}\gets\mathit{tail}+v$ \textbf{else} $w'[c]\gets w'[c]+v$
  \EndFor
  \State $w\gets w'$
\EndFor
\State \Return $\mathit{tail}$
\end{algorithmic}
\caption{흡수 경계가 있는 두 표본 Kolmogorov--Smirnov 분포}
\label{alg:ks-absorb}
\end{algorithm}

상태 $c$는 $p$걸음 뒤에 경계 안에 남아 있어야 하므로, 양쪽 검정에서는 $|c(\alpha+\beta)-p\beta|<h$를 만족하는 띠 안에만 있다. 띠의 폭은 $2h/(\alpha+\beta)+1$ 이하이므로 일은 $N$에 띠의 폭을 곱한 정도이다. 또 무게가 아래로 넘쳐 $0$이 된 상태는 띠에서 빼도 결과에 영향이 없으므로, 실제로 살아 있는 상태는 더 적다.

동점이 있으면 같은 값 $t$개가 한 덩어리로 들어오고, 경험 분포 함수는 덩어리의 끝에서만 관측된다. 덩어리가 첫째 표본의 값을 $k$개 가질 확률은 남은 $R$자리에 첫째 표본이 $r=m-c$개 남은 상황의 초기하 확률 $\binom tk\binom{R-t}{r-k}/\binom Rr$이므로, \cref{alg:ks-absorb}의 두 갈래를 이 확률로 갈라지는 여러 갈래로 바꾸면 된다. 이것은 Mann--Whitney에서처럼 관측한 동점 구조에 조건부인 분포이다.

\begin{instardust}[\code{csrc/htest/ks2.hpp}]
\code{tail}은 \cref{alg:ks-absorb}를 double-double로 한 것이며, 크기가 다른 두 표본의 차를 정수 $g=i\alpha-j\beta$로 다루고 경계를 정수끼리 비교한다. 동점이 없을 때는 새 상태 $c$가 $w(c-1)$에서 첫째 표본으로 받는 기여와 $w(c)$에서 둘째 표본으로 받는 기여를 한꺼번에 계산하고, 경계 안의 연속 구간과 밖으로 흡수되는 양 끝을 나누어 처리한다. 동점이 있으면 덩어리마다 \cref{sec:counting}의 \code{hypergeometric}으로 갈래의 확률을 만든다. 크기가 같고 동점이 없으면 \code{equal\_tail}이 \cref{prop:two-sided-reflection}의 닫힌 꼴을 비의 곱으로 계산한다. scipy는 경로의 수를 \code{double}로 세기 때문에 $\binom Nm$이 약 $10^{308}$을 넘으면 정확 계산을 포기하고 점근식으로 넘어가지만, 확률로 쌓는 이 방법에는 그런 한계가 없다.
\end{instardust}

\section{맺음}

이 장에서 본 정확 분포들은 모두 같은 틀을 따른다. 귀무가설이 관측의 배치를 균등하게 만들면 p값은 개수의 몫이 되고, 통계량이 독립적인 선택들의 합이면 그 분포의 생성함수는 인수들의 곱이 된다. Mann--Whitney는 $q$-이항계수, 부호 순위는 $\prod(1+q^i)$, Kendall은 $\prod[j]_q$, 동점이 있는 경우는 두 변수 곱이었고, 격자 경로는 반사 원리와 흡수 경계로 다루었다. 계산의 측면에서는 세 가지 원칙이 거듭 나왔다. p값에 필요한 꼬리의 차수에서 자르고(\cref{prop:fps-truncate}), 대칭이 있으면 가까운 쪽 꼬리만 세며, 정수의 범위를 넘으면 뺄셈 없는 양수의 합으로 확률을 쌓는다. 마지막 원칙이 왜 상대오차를 묶어 두는지는 2부에서 반올림 오차를 정식으로 해석하면서 증명한다.

조합론의 생성함수 방법은 \cite{wilf2006gf}가 가장 친절한 입문이고, $q$-이항계수와 분할은 \cite{andrews1998partitions}, 세기 문제 전반은 \cite{stanley2012ec1}과 \cite{flajolet2009analytic}가 표준적인 참고서이다. 순위 검정의 이론은 \cite{hajek1999rank}에 정리되어 있다.

\chapterreferences
